\chapter{射影簇}\label{chap:projective-varieties}

\noindent\textbf{第 \thechapter{} 章学习目标}
\begin{itemize}
  \item 从不同角度定义并使用射影空间.
  \item 理解多项式在射影空间一点处取值为零的含义, 以及齐次性在射影消没中的作用.
  \item 在射影情形定义并使用 $\mathcal{V}$ 和 $\mathcal{I}$ 算子, 理解它们之间的关系.
  \item 计算射影簇的仿射限制和仿射簇的射影闭包, 在仿射与射影两种情形间转换.
  \item 陈述射影零点定理, 并通过仿射锥把它归结为仿射零点定理.
\end{itemize}

到这里为止, 本书研究过的所有簇都位于仿射空间 $\mathbb{A}^{n}$ 中. 不过, 还有一个更大的环境空间, 在其中\textquotedbl{}簇\textquotedbl{}的概念也有意义. 这个空间称为\emph{射影空间}, 记为 $\mathbb{P}^{n}$. 本章的目标是定义 $\mathbb{P}^{n}$, 以及在 $\mathbb{P}^{n}$ 中作为多项式零点集得到的\emph{射影簇}.

这种推广的动机是, 要在关于仿射簇的陈述不可避免地出现例外时, 给出没有这些例外的统一陈述. 一个典型例子是: 在 $\mathbb{A}^{2}$ 中, 任意两条直线都相交, 平行线除外. 可以把射影空间 $\mathbb{P}^{2}$ 看成在 $\mathbb{A}^{2}$ 中添加了\textquotedbl{}无穷远点\textquotedbl{}, 使 $\mathbb{A}^{2}$ 中每条直线都经过由其斜率决定的一个特定无穷远点. 添加这些点后, $\mathbb{P}^{2}$ 中每两条直线都相交, 无一例外. 这是一个优美结论的特例, 这个结论称为 Bézout 定理: 在 $\mathbb{P}^{2}$ 中, 由次数为 $r$ 和 $s$ 的多项式定义的两条曲线, 若以适当方式计数, 则相交于 $r\cdot s$ 个点. 虽然 $\mathbb{A}^{2}$ 中相应的陈述往往为真, 但容易找到例外. 例如, 抛物线 $\mathcal{V}(y-x^{2})$ 和竖直直线 $\mathcal{V}(x)$ 只相交于一点. 从 $\mathbb{P}^{2}$ 的角度看, 这个例外仍是因为在\textquotedbl{}无穷远处\textquotedbl{}还有一个交点, 只关注 $\mathbb{A}^{2}$ 时它是看不到的.

这些平面几何中的观察最终使代数几何学家认识到, 研究多项式方程组的解时, 最自然的环境空间是 $\mathbb{P}^{n}$, 而非 $\mathbb{A}^{n}$. 射影空间的定义初读时可能难以消化, 但它为代数几何研究带来的简洁与统一, 值得我们付出这番努力.

\par\smallskip
{\footnotesize
\noindent\textcopyright{} The Author(s) 2025\par
\noindent E. Clader, D. Ross, \emph{Beginning in Algebraic Geometry}, Undergraduate Texts in Mathematics,
\href{https://doi.org/10.1007/978-3-031-88819-9_10}{\nolinkurl{https://doi.org/10.1007/978-3-031-88819-9_10}}.\par}
